\chapter{EKSPERIMENTALNI REZULTATI I RASPRAVA}
\label{ch:rezultati}

Poglavlja \ref{ch:klasicni} i \ref{ch:rl} iznose rezultate svakog pristupa na njegovim
vlastitim uvjetima. Ovdje se ta dva pristupa međusobno uspoređuju. Najprije se utvrđuje što
se uopće smije uspoređivati, zatim izlaže usporedba po zajedničkom kriteriju uspjeha,
kontaktnim silama i načinima otkazivanja, i naposljetku raspravlja o tome što razlike znače.

\section{Protokol usporedbe}
\label{sec:metrike}

Oba su pristupa mjerena na sceni sa zidovima, uz otpor vrata izvučen iz istog raspona
(tablica \ref{tab:randomizacija}), ali postupci mjerenja razlikuju se u trima pogledima koji
određuju koje su tvrdnje opravdane.

Klasični pristup izvodi cijeli lanac, od prilaska i hvata preko
otvaranja do prolaska kroz vrata. Pristup temeljen na učenju rješava isključivo fazu otvaranja i
kreće od stanja u kojem prihvatnica već drži kvaku (potpoglavlje \ref{sec:rl-opseg}). Usporedba
se zato odnosi na otvaranje.

Klasični je pristup mjeren u deset zasebnih pokušaja po tipu vrata,
naučena politika u nekoliko stotina epizoda.

Pragovi se razlikuju: klasični pristup cilja 50\textdegree{} odnosno
0,90 m, a naučena politika 100\textdegree{} odnosno 0,60 m. Mjerenja obaju pristupa zato bilježe
ostvareni otvor kao kontinuiranu veličinu, pa se stopa uspjeha računa za bilo koji prag
naknadno, upravo zato da usporedba ne ovisi o tome čiji je kriterij odabran.

Pri determinističkom izvođenju naučene politike dio epizoda ne ostvaruje gibanje platforme, pa
manipulator doseže granicu vlastitog radnog prostora i otvaranje staje na oko 0,45 m odnosno
12\textdegree{}. Takve su epizode izdvojene iz statistike, pa analiza obuhvaća 305 od 562
epizode na zakretnim te 286 od 581 na kliznim vratima; uzrok je opisan u potpoglavlju
\ref{sec:rl-limits}. Navedene brojke time opisuju sposobnost upravljanja, a ne pouzdanost
izvođenja u cjelini.

\section{Uspješnost otvaranja}
\label{sec:usporedba}

Tablica \ref{tab:usporedba-uspjeh} uspoređuje oba pristupa po pragu koji koristi klasični
pristup i po pragu koji koristi naučena politika.

\begin{table}[H]
    \centering
    \caption{Stopa uspjeha po zajedničkim pragovima}
    \label{tab:usporedba-uspjeh}
    \begin{tabular}{@{}llcc@{}}
        \toprule
        Vrata    & Prag            & Klasični pristup & Pristup učenjem \\
        \midrule
        zakretna & 50\textdegree{} & 0,80             & 1,00            \\
        zakretna & 10\textdegree{} & 0,00             & 0,96            \\
        klizna   & 0,60 m          & 1,00             & 1,00            \\
        klizna   & 0,90 m          & 1,00             & 1,00            \\
        \bottomrule
    \end{tabular}
\end{table}

Kod kliznih vrata oba pristupa dosežu puni cilj od 0,90 m u svakom pokušaju. Zadatak je
za oba rješiv: ograničenje je pravac poznat iz percepcije kod klasičnog pristupa, odnosno pravac
koji politika otkriva iz sile i pomaka, a smjer sile poklapa se sa smjerom gibanja. Razlika
među pristupima ondje se ne očituje u uspješnosti.

Kod zakretnih vrata razlika je najizraženija u cijelom radu. Naučena politika otvara
vrata preko 100\textdegree{} u 96 \% epizoda, dok klasični pristup taj kut ne doseže
nijednom, ne zato što ga ne cilja, nego zato što mu je fizička granica oko
57\textdegree{}, iznad koje kvaka klizne iz prihvatnice (potpoglavlje \ref{sec:klasicni-limits}).

Razlika nije samo u dosegnutom pragu nego i u obliku raspodjele. Klasični pristup staje na
52,12\textdegree{} $\pm$ 0,42\textdegree{}, dakle vrlo ponovljivo i skoro točno ondje gdje mu je cilj
postavljen. Naučena politika ima medijan otvora na 122\textdegree{}, a gornji kvartil na
134,1\textdegree{}, praktički na mehaničkoj granici zgloba od 135\textdegree{}. Politika
vrata redovito otvara do kraja hoda. Ista se slika ponavlja i kod kliznih vrata, gdje
medijan otvora leži na punom hodu od 1,00 m.

\section{Kontaktne sile}
\label{sec:usporedba-sile}

Izmjerene sile klasičnog pristupa navedene su u tablici \ref{tab:usporedba-sile}. Naučena
politika prosječnu silu ne izvještava izravno, ali je njezina raspodjela ograničena sigurnosnim
prekidom na 400 N: udio epizoda prekinutih prekoračenjem sile iznosi 0,6 \% kod zakretnih i
2,8 \% kod kliznih vrata (tablica \ref{tab:rl-rezultati}). Politika dakle u više od 97 \%
epizoda ostaje ispod te granice.

\begin{table}[H]
    \centering
    \caption{Kontaktne sile klasičnog pristupa, prosjek deset pokušaja}
    \label{tab:usporedba-sile}
    \begin{tabular}{@{}lcc@{}}
        \toprule
        Vrata    & Sila, prosječna & Sila, najveća \\
        \midrule
        zakretna & 348 N           & 625 N         \\
        klizna   & 258 N           & 661 N         \\
        \bottomrule
    \end{tabular}
\end{table}

Klasični pristup redovito prelazi 400 N, s vršnom silom od 625 N, dok naučena politika tu granicu
prekoračuje u manje od 1 \% posto epizoda.
Kod kliznih vrata klasični pristup generalno ostaje ispod granice, ali za otići i preko (vrh 661 N),
pa je naučena politika i tu bolja.

Taj je nalaz izravna potvrda teze koju rad gradi od poglavlja \ref{ch:klasicni}. Klasični pristup
vodi vrh alata krutim pozicijskim upravljanjem, pa svako odstupanje između procijenjene i stvarne
geometrije luka postaje sila. Naučena politika krutost regulatora bira po osi i u svakom koraku,
pa silu smanjuje popuštanjem u smjeru u kojem se bori s ograničenjem.

\section{Načini otkazivanja}
\label{sec:usporedba-otkazi}

Kod klasičnog pristupa oba neuspjeha na zakretnim vratima dogodila su se u fazi hvata,
prije nego je otvaranje uopće počelo. Kad otvaranje jednom krene, uspije svaki put i s vrlo
malim rasipanjem. Slabija karika klasičnog pristupa kod zakretnih vrata time nije upravljanje
nego percepcija i hvat.

Naučena politika otkazuje tijekom same interakcije: gubi hvat u 7,2 \% epizoda na zakretnim
vratima i dodiruje krilo u 6,1 \%. Kod kliznih su vrata oba gotovo nula. Dodir zida je u oba
slučaja zanemariv, što pokazuje da je politika naučila pouzdano izbjegavati prepreke.

\section{Rasprava}
\label{sec:rasprava}

\subsection{Različiti ulazi u isti problem}

Razlika u uspješnosti ne smije se čitati kao razlika u kvaliteti upravljanja, jer dva pristupa ne
rješavaju problem s istim informacijama.

Klasični pristup zna gdje je os šarke. Izračunava je iz poze oznake i poznate širine
krila, kao fiksni vektor u koordinatnom sustavu oznake (izraz \eqref{eq:axis}), uz izmjerenu
pogrešku od 5 do 8 mm. Naučena politika taj podatak nema ni u kojem obliku: geometriju
ograničenja mora otkriti iz vlastitih opažanja sile i pomaka, a kut vrata nikada ne ulazi u
njezino opažanje. Ista asimetrija vrijedi i za smjer klizanja kod kliznih vrata.

Asimetrija ide i u drugom smjeru. Naučena politika koordinira platformu i manipulator jednom,
naučenom strategijom (potpoglavlje \ref{sec:baza-uci}), dok klasični pristup taj problem rješava
unaprijed propisanim slijedom naizmjeničnih faza. Kod zakretnih vrata upravo ta propisana shema
postaje ograničenje: izmjereno je da primicanje platforme manipulatoru ne vraća doseg, jer kvaka
putuje zajedno s platformom, pa doseg kroz cijelo pokretanje ostaje između 0,78 i 0,87 m.
Politika koja koordinaciju uči nije vezana tom pretpostavkom i pronalazi obrazac gibanja koji
propisani slijed ne sadrži.

\subsection{Zašto se razlika pojavljuje samo na zakretnim vratima}

Izostanak razlike na kliznim i velika razlika na zakretnim vratima nisu slučajni, nego slijede iz
toga koliko svaki zadatak traži prilagodbu krutosti tijekom interakcije.

Klizna su vrata geometrijski jednostavna: ograničenje je pravac, smjer sile poklapa se sa smjerom
gibanja, a jedina poteškoća je duljina hoda, koju klasični pristup rješava ponavljanjem prolaza.
Ondje eksplicitno poznavanje geometrije daje sve što treba, pa dodatna sloboda u odabiru krutosti
ne donosi mjerljivu prednost.

Zakretna vrata traže gibanje po luku, uz stalno mijenjanje odnosa između platforme, manipulatora
i kvake. Ondje kruto praćenje unaprijed izračunate putanje stvara silu na svako odstupanje, a
uhvaćena kvaka pritom klizne. Naučena politika, koja krutost mijenja po osi i koordinaciju uči,
prolazi kroz područje na koje klasični pristup nailazi kao na tvrdu granicu.

Zaključak je time uvjetovan, a ne općenit: eksplicitno poznavanje geometrije dostatno je kad je
geometrija jednostavna i dobro procijenjena, a gubi prednost čim zadatak zahtijeva prilagodbu
krutosti tijekom same interakcije.

\section{Ograničenja usporedbe}
\label{sec:usporedba-limits}

\begin{itemize}
    \item Usporedba obuhvaća isključivo fazu otvaranja. Prilazak, hvat i prolazak kroz vrata
          klasični pristup izvodi, a naučena politika ne, pa se njihova ukupna sposobnost
          rješavanja zadatka ovim mjerenjima ne uspoređuje.

    \item Navedene stope uspjeha naučene politike uvjetne su na način opisan u potpoglavlju
          \ref{sec:metrike} i opisuju sposobnost upravljanja, a ne pouzdanost izvođenja u
          cjelini.

    \item Stopa uspjeha nije istovrsna veličina na dvjema stranama. Kod klasičnog pristupa riječ
          je o binarnom ishodu po pokušaju, nad deset pokušaja; kod naučene politike o udjelu
          nad nekoliko stotina epizoda. Razlika između 0,80 i 1,00 na uzorku od deset pokušaja
          nije pouzdana; razlika između 0,00 i 0,96 jest.

    \item Naučena politika mjerena je uz nasumičan kut prilaska, kojega u klasičnim pokušajima
          nema, pa raspon uvjeta nije posve izjednačen.

    \item Trajanje izvođenja nije uspoređivano. Epizoda naučene politike traje najviše deset
          sekundi i najčešće istječe bez prekida, jer uspjeh ne prekida epizodu, dok klasični
          pristup traje dok zadatak ne dovrši, prosječno 40 s na zakretnim i 28 s na kliznim
          vratima. Te dvije veličine ne mjere isto.

    \item Obje su izvedbe mjerene isključivo u simulaciji. Zaključci se odnose na model, a
          prenosivost na stvarni robot ograničena je pojednostavljenjima navedenima u poglavljima
          \ref{ch:sustav} i \ref{ch:percepcija}.
\end{itemize}
